% =====================================================================================
% Chapter 12 · Redesign: interrupted time series
%
% Built from notebooks/10_redesign_its.ipynb. NOT ONE NUMBER IN THIS FILE IS TYPED BY
% HAND: every figure in the prose is a macro from build/macros.tex, generated by
% cmp.report from the executed notebook. The only literals below are (i) the constants
% of the data-generating model in §12.2, which are *definitions* rather than results,
% and (ii) the decision-rule thresholds (0.9 / 0.5), which are *conventions* stated in
% the text, not measurements.
% =====================================================================================
% \pp is typography, not a number: percentage points, set like the preamble's \pc.
% (Defined here rather than in the shared preamble, which this chapter does not own.)
\providecommand{\pp}[1]{\mbox{#1\,pp}}

\chapter{Redesign: interrupted time series}
\label{ch:its}

\begin{quote}\small
\emph{A homepage redesign shipped on one day, to every visitor, everywhere. There is no control
group, no untreated geography, and there never will be: the only comparison the world offers is
the site's own past. Interrupted time series fits that past --- level, trend, weekly cycle ---
extrapolates it across the \nbTenNPost{} days that follow, and reads the gap, which comes to
\pp{\nbTenBayesAvg} against a planted \pp{+\nbTenTrueLift} where the dashboard's before-and-after
number said \pp{\nbTenNaive}. On a \eur{\nbTenCost} build the verdict is \nbTenPersVerdictnodecay{}.
And then an unrelated \pp{+\nbTenShock} shock is planted on the redesign date, and both arms ---
the segmented regression with its Newey--West interval and the posterior with its credible band ---
report \pp{\nbTenShockCls} and \pp{\nbTenShockBayes}, each inside a tight interval that
\emph{excludes} the truth, while every in-sample diagnostic stays green. Nothing in the apparatus
bought protection: the credible interval and the confidence interval both price sampling noise, and
the ITS problem is not sampling noise. \Cref{sec:its:wrong} is the chapter.}
\end{quote}

% =====================================================================================
\section{The decision}
\label{sec:its:decision}

The homepage was redesigned on a known day. It shipped to all traffic at once --- no holdout
bucket, no phased rollout, no untreated market --- and conversion has been higher ever since. The
proposal on the table is to keep it, at an amortized build-and-maintenance cost of
\eur{\nbTenCost} a year; the alternative is to roll it back.

The question is not whether conversion rose. It did. The question is whether the \emph{redesign}
raised it. The site was already drifting upward before the change --- steadily, unremarkably,
for reasons nobody has ever fully attributed --- and a before-and-after comparison of means books
that drift as redesign lift. On this series it books \pp{\nbTenNaiveTrendBias} of it, which is
not a rounding error on an effect of \pp{+\nbTenTrueLift}: it is a third of the answer, and it
points the wrong way.

What makes the problem hard is that every escape route used elsewhere in this book is closed.
There is no coin flip, so the randomized comparison behind a geo experiment is
unavailable. There is no untreated geography, so synthetic control has no donor
pool. There is no staggered rollout, so difference-in-differences has no never-treated cohort to
lean on. Everyone was treated, on the same day, at the same moment. \textbf{The only untreated
comparison in existence is the site's own history.}

Interrupted time series (ITS) is the method that uses it \citep{bernal2017}. Fit the pre-change
series --- its level, its trend, its weekly rhythm --- extrapolate that fit forward past the
change date, and call the extrapolation the counterfactual: what conversion \emph{would} have done
had nothing shipped. The gap between what actually happened and that projection is the estimated
effect. Stated that plainly, the method's weakness is already visible, and this chapter's job is
to measure it rather than to apologise for it: the counterfactual is a \emph{forecast}, made into
a window where nothing can check it.

A wrong answer costs in both directions. Roll back a redesign that works and the firm walks away
from \eur{\nbTenPersValuenodecay} a year of conversion value. Keep one that does not, and it has
paid \eur{\nbTenCost} for a coincidence and will build the next one on the same false lesson. The
chapter estimates the lift classically, re-estimates it as a posterior, puts the two on the same
estimand --- and then breaks the assumption that both of them rest on, to show that neither
apparatus noticed.

% =====================================================================================
\section{The data-generating model}
\label{sec:its:dgm}

A causal method cannot be graded on real data, because on real data the counterfactual is missing
--- that is the entire problem, and in this chapter it is the problem in its purest form. So the
method is graded against a world whose answer is known by construction.

The simulator produces \nbTenNDays{} days of a daily conversion rate,
$t = 0, \dots, \nbTenNDays-1$, with the redesign landing on day $t^{*} = \nbTenRedesignDay$ for
everyone at once. Three ingredients, and only the third of them is the effect:
%
\begin{equation}
  Y_t \;=\; \underbrace{0.10 + 0.00015\,t}_{\text{level {+} drift}}
        \;+\; \underbrace{0.01 \sin\!\Big(\tfrac{2\pi t}{7}\Big)}_{\text{weekly cycle}}
        \;+\; \underbrace{0.03\;\mathbf 1[\,t \ge t^{*}\,]}_{\text{the redesign}}
        \;+\; \varepsilon_t,
  \qquad \varepsilon_t \sim \mathcal N\!\left(0,\ 0.006^{2}\right).
  \label{eq:its:dgp}
\end{equation}
%
The planted effect is a \emph{pure level step} of \pp{+\nbTenTrueLift} --- about
\pc{\nbTenTrueLiftRel} in relative terms on a base of \pc{\nbTenBaseRate} --- switched on by the
indicator $\mathbf 1[t \ge t^{*}]$ and constant thereafter. It is the grade every estimator in
this chapter is marked against, and the last thing in the chapter that is known for certain.

The drift is the confounder, and its size is worth computing before any estimator runs. At
\pp{\nbTenDriftDay} per day it compounds to \pp{+\nbTenDriftPre} over the \nbTenNPre-day
pre-period. A before-and-after difference of means compares a post-window centred on day
$\approx 150$ with a pre-window centred on day $\approx 60$: those centres sit
\nbTenNaiveGapDays{} days of drift apart, so \emph{arithmetic alone} predicts that roughly
\pp{\nbTenNaiveTrendBias} of pure pre-existing trend will be booked as redesign lift. The
dashboard number duly comes in at \pp{\nbTenNaive} against a truth of \pp{+\nbTenTrueLift}
(\cref{tab:its:naive}) --- and the entire methodological content of ITS is the subtraction that
removes it.

The weekly cycle is not decoration either. It is the one genuinely autocorrelated feature of the
series, and \cref{sec:its:valid} shows what happens to the diagnostics when a model forgets it.

\input{build/tables/nb10_readings}

% =====================================================================================
\section{Identification}
\label{sec:its:ident}

\subsection*{The estimand}

In the potential-outcome notation of \citet{rubin1974} and \citet{imbens2015}, write $Y_t(1)$ for
the conversion observed on day $t$ with the redesign live and $Y_t(0)$ for the conversion that
would have been observed had nothing shipped. After $t^{*}$ only the first is observed; before
$t^{*}$ only the second. The daily effect is their contrast, and the object the business is buying
is its average over the post-redesign window:
%
\begin{equation}
  \Delta_t \;=\; Y_t(1) - Y_t(0), \qquad t \ge t^{*},
  \qquad\qquad
  \bar\Delta \;=\; \frac{1}{T - t^{*}} \sum_{t \ge t^{*}} \Delta_t ,
  \label{eq:its:estimand}
\end{equation}
%
with $T = \nbTenNDays$ and hence \nbTenNPost{} post-redesign days. Its cumulative sibling,
%
\begin{equation}
  C_T \;=\; \sum_{t \ge t^{*}} \Delta_t \;=\; (T - t^{*})\,\bar\Delta ,
  \label{eq:its:cum}
\end{equation}
%
is what \cref{sec:its:euro} turns into realized euros. \Cref{eq:its:cum} is a \emph{linear} map of
\cref{eq:its:estimand}, and that innocuous observation is the hinge of \cref{sec:its:compare}.

Everything in this chapter --- both estimators, both intervals, and the euro decision --- targets
$\bar\Delta$. Keeping the estimand fixed is what makes the comparison a comparison rather than a
change of subject.

\subsection*{The assumption that has to do all the work}

$Y_t(0)$ is never observed after $t^{*}$. It must be \emph{reconstructed}, and the only material
available is the pre-period. So the identifying claim is a claim about a structure surviving a
date:
%
\begin{equation}
  \E\big[Y_t(0)\big] \;=\; \beta_0 + \beta_1 t + \beta_2 \sin\!\Big(\tfrac{2\pi t}{7}\Big)
                          + \beta_3 \cos\!\Big(\tfrac{2\pi t}{7}\Big)
  \qquad\text{for \emph{all} } t, \text{ not merely } t < t^{*}.
  \label{eq:its:ident}
\end{equation}
%
Read \cref{eq:its:ident} carefully: it says that \textbf{no other term appears at $t^{*}$}. Not
a price change, not a competitor's outage, not a coincident campaign, not a tracking migration,
not a seasonality break. If any of them lands on the redesign date, it enters the equation the
model does not have, and its effect is attributed to the term the model does have.

Five assumptions make \cref{eq:its:ident} respectable. They are stated here, and each is either
tested later in the chapter or explicitly declared untestable --- an assumption with no
accompanying test is a wish, and \cref{as:its:shock} is the wish.

\begin{assumption}[No concurrent shock]\label{as:its:shock}
Nothing other than the redesign changes the data-generating process at $t^{*}$. Formally, the
missing-term statement of \cref{eq:its:ident}. \textbf{This assumption is untestable from the
series.} A shock arriving with the redesign is observationally identical to a larger redesign
effect: no statistic computed from $\{Y_t\}$ can separate them, because they imply the same
distribution of the same data. It is discharged only from outside the data --- by auditing the
launch calendar --- and \cref{sec:its:wrong} plants a violation on purpose to show what its
failure costs and, more importantly, what it looks like from the inside, which is: nothing.
\end{assumption}

\begin{assumption}[Correct pre-period structure]\label{as:its:form}
The functional form in \cref{eq:its:ident} --- level, linear trend, one weekly Fourier pair --- is
the right description of the untreated series. A missing term does not merely add noise; it lands
in the residual and distorts the extrapolation. This one \emph{is} testable, in two ways: the
residual autocorrelation of the pre-period fit (\cref{sec:its:valid}, where dropping the seasonal
pair makes the lag-7 bar jump from \nbTenAcfSeven{} to \nbTenAcfBadSeven), and a pre-period
backtest that grades the forecast where grading is still possible.
\end{assumption}

\begin{assumption}[No anticipation]\label{as:its:antic}
Nothing shifts conversion \emph{before} $t^{*}$ --- no leaked beta, no staged ramp, no mis-dated
launch. Tested by placebo-in-time: refitting with a fake redesign at day \nbTenFakeDay{} returns
\pp{\nbTenPlacebo}, and a sweep over fake dates (\nbTenPlaceboDates) returns at most
\pp{\nbTenPlaceboMax} in absolute value.
\end{assumption}

\begin{assumption}[Exogenous intervention date]\label{as:its:date}
The launch date was not chosen \emph{because} of the metric. Redesigns rarely ship on random days;
they ship because conversion looked bad. If a transient dip triggered the launch, the pre-period
fit inherits the dip and the natural rebound is booked as lift --- regression to the mean wearing
a causal costume. The placebo-date sweep of \cref{sec:its:valid} probes it (a real dip would make
nearby fake dates light up), but the decisive check is institutional: ask what triggered the date.
\end{assumption}

\begin{assumption}[Weak dependence of the errors]\label{as:its:dep}
The residuals of \cref{eq:its:ident} are not so persistent that an interval built on them is
meaningless. This is not an identification assumption --- it does not affect the point estimate ---
but it governs every interval in the chapter, and it is the one place where the classical arm's
choice of covariance estimator earns or loses its keep. Tested by the residual ACF; priced, in
\cref{sec:its:classical}, by rebuilding the series with AR(1) errors and measuring what each
standard error's interval actually covers.
\end{assumption}

The asymmetry between \cref{as:its:shock} and the rest is the intellectual content of this
chapter. Four of the five can be interrogated with data. The one that carries the causal claim
cannot.

% =====================================================================================
\section{The classical baseline}
\label{sec:its:classical}

Before any sampler, the discipline of this book is to ask what a competent analyst produces
\emph{without} one. Here that analyst reaches for the workhorse of the applied ITS literature ---
\textbf{segmented regression} \citep{wagner2002, bernal2017}, one ordinary least-squares fit on the
whole series in which the intervention enters as \emph{two} terms:
%
\begin{equation}
  Y_t \;=\; \beta_0 + \beta_1 t
        \;+\; \underbrace{\delta\,\mathbf 1[t \ge t^{*}]}_{\text{level shift}}
        \;+\; \underbrace{\gamma\,(t - t^{*})^{+}}_{\text{slope change}}
        \;+\; \beta_2 \sin\!\Big(\tfrac{2\pi t}{7}\Big)
        \;+\; \beta_3 \cos\!\Big(\tfrac{2\pi t}{7}\Big)
        \;+\; \varepsilon_t ,
  \label{eq:its:seg}
\end{equation}
%
with $(t - t^{*})^{+} \equiv \max(t - t^{*},\, 0)$. Read the coefficients one by one. $\beta_1$ is
the drift the dashboard mistook for lift. $\delta$ is the jump on the redesign day. $\gamma$ is
the change of \emph{trajectory} afterwards --- zero if the redesign is a pure step. And
$\beta_2, \beta_3$ absorb the weekly cycle.

Setting $\delta = \gamma = 0$ and evaluating the fitted line for $t \ge t^{*}$ gives
$\widehat{Y_t(0)}$: the same counterfactual, extrapolated out of the same pre-period structure into
the same uncheckable window. \textbf{The classical arm is not a safer analysis that dispenses with
the leap of faith; it is the same leap, taken with lighter luggage.} Nothing in \cref{eq:its:seg}
knows whether the pre-period shape survived day $t^{*}$, and \cref{sec:its:wrong} refits this very
equation on a series where it did not.

\subsection*{The point estimate}

Fitted to all \nbTenNDays{} days, \cref{eq:its:seg} returns a level shift of
$\hat\delta = \pp{\nbTenClLevel}$, with a 90\,\% interval of
$[\pp{\nbTenClLevelLo},\ \pp{\nbTenClLevelHi}]$
and $R^2 = \nbTenClRsq$. Against the planted \pp{+\nbTenTrueLift} that is an error of
\pp{\nbTenClLevelErr}, and the truth falls inside the interval. The trend model earned its keep:
it stripped out \pp{\nbTenClDriftRemoved} of drift that the naive number had reported as lift
(\cref{tab:its:naive}).

The second coefficient is worth as much as the first, and it is one the Bayesian section of this
chapter can only eyeball. The slope change comes back at
$\hat\gamma = \nbTenClSlope$\,pp/day, 90\,\% interval
$[\nbTenClSlopeLo,\ \nbTenClSlopeHi]$ --- straddling zero. The redesign is a \textbf{step}, not a
trajectory change, matching \cref{eq:its:dgp}. That matters commercially: a single average lift is
a fair summary of the whole post-window and will not drift as the reporting window lengthens. Had
$\gamma$ been significant, the honest deliverable would have been two numbers rather than one, and
every annualisation in \cref{sec:its:euro} would have needed rebuilding.

To report the summary the posterior reports --- the mean post-period gap of
\cref{eq:its:estimand} --- note that in the parameters of \cref{eq:its:seg},
%
\begin{equation}
  \bar\Delta \;=\; \delta \;+\; \gamma \cdot \overline{(t - t^{*})^{+}}
             \;=\; \delta \;+\; 29.5\,\gamma ,
  \label{eq:its:recentre}
\end{equation}
%
the mean of $(t - t^{*})^{+}$ over the \nbTenNPost{} post days being $29.5$. Rather than
delta-method that combination, the slope term is simply \emph{re-centred} --- subtract $29.5$ from
$(t - t^{*})^{+}$ on the post days --- after which the coefficient on the post indicator \emph{is}
$\bar\Delta$ and the fit hands over its standard error directly. Same regression, reparameterised
so that the estimand is a coefficient. It gives $\pp{\nbTenClAvg}$, 90\,\% interval
$[\pp{\nbTenClAvgLo},\ \pp{\nbTenClAvgHi}]$. That is the number \cref{sec:its:compare} lays beside
the posterior.

\subsection*{The covariance choice, argued and then demonstrated}

A point estimate without a defensible standard error is half an answer, and here the default
standard error is the indefensible half. Daily conversion is a \emph{time series}: today's residual
tends to resemble yesterday's, because an unmodelled promotion runs for a week, a tracking glitch
lasts three days, a competitor's outage bleeds across a weekend. The iid formula prices
\nbTenNDays{} days as \nbTenNDays{} independent facts. The classical repair is the
\textbf{Newey--West HAC} standard error \citep{newey1987} --- heteroskedasticity- \emph{and}
autocorrelation-consistent --- which allows the residuals to correlate up to a chosen lag and
charges the estimate for that correlation.

The lag is a choice, and it is made twice over. Substantively, \textbf{lag~7 is one full weekly
cycle}, the natural dependence horizon of a daily web series. Mechanically, the common rule of
thumb $\lfloor 4(n/100)^{2/9} \rfloor$ returns \nbTenHacLagRule{} at $n = \nbTenNDays$, so the
choice made here is the \emph{more} conservative one, and the estimate is in any case insensitive
to it: the standard error on $\bar\Delta$ moves from \nbTenHacSeLagfour{} at lag~4 to
\nbTenHacSeLagseven{} at lag~7, \nbTenHacSeLagfourteen{} at lag~14 and
\nbTenHacSeLagtwentyone{} at lag~21 --- a decision that changes no decision.

The textbook slogan is that HAC widens the interval. On this series it does not, and the honest
thing to do with a slogan that fails is to say why rather than to bury the measurement. The HAC
standard error on $\bar\Delta$ is \nbTenClSe{} against the iid \nbTenClSeIid{} --- a ratio of
$\nbTenClSeRatio\times$, which is a rounding error. This is the argument \emph{working}, on data
that happens not to need it: the noise of \cref{eq:its:dgp} is iid by construction, and the one
genuinely autocorrelated feature of the series --- the weekly cycle --- is \emph{already in the
formula}. Newey--West goes looking for leftover serial correlation, finds none, and returns the iid
answer plus its own estimation noise.

So why pay for it? Because of \cref{tab:its:hac}, which isolates the mechanism instead of asserting
it. The same series is rebuilt with AR(1) errors of persistence $\rho$ --- same level, same trend,
same weekly cycle, same planted \pp{+\nbTenTrueLift} step, same marginal noise, \emph{only} the
serial dependence changes --- and each covariance estimator is asked the only question that
ultimately matters about a standard error: \textbf{does its 90\,\% interval cover the truth
90\,\% of the time?}

\input{build/tables/nb10_hac}
\input{build/floats/nb10_hacfig}

Read \cref{tab:its:hac} across, not down. The iid standard error \emph{does not move} as $\rho$
rises. This is not underestimation so much as blindness: serial correlation is not a quantity that
appears anywhere in the iid formula, so it cannot be perceived, and what gives way instead is the
one thing the analyst cares about. The iid interval's coverage falls from
\pc{\nbTenArCovIidZero} to \pc{\nbTenArCovIidEight} --- at $\rho = 0.8$ a nominal ninety is closer
to a coin flip than to its label --- while the HAC interval holds \pc{\nbTenArCovHacZero} to
\pc{\nbTenArCovHacEight}.

Two lessons, and the second is the one practitioners skip. First: \textbf{HAC is insurance that
costs almost nothing when the residuals are clean and is the difference between a 90\,\% interval
and a coin flip when they are not} --- and one cannot know which world one is in without computing
it. Real daily conversion data is never the clean world. Second: HAC is a \emph{repair, not a
cure}. Its own coverage erodes at high persistence, to \pc{\nbTenArCovHacEight}, because it must
estimate the autocorrelation from the same short series. With strongly persistent residuals,
\emph{structure beats a sandwich}: model the dependence (an AR term, or the state-space treatment
of \citealt{brodersen2015}) rather than patch it after the fact.

\subsection*{What the classical arm cannot say}

An OLS fit with five regressors and one honest sandwich has recovered the truth, without a
likelihood, a prior, or a sampler. That is the baseline the next section must justify itself
against. But note the boundary of what it produced. The interval
$[\pp{\nbTenClAvgLo},\ \pp{\nbTenClAvgHi}]$ is \textbf{not} a 90\,\% probability that the lift
lies inside it: it is a property of the \emph{procedure}, and this particular interval either
covers the truth or does not. \Cref{sec:its:euro} must answer \emph{``what is the probability the
redesign's annual value clears its \eur{\nbTenCost} cost?''} --- a probability about the
\emph{effect}. The classical apparatus has none to give, not because its arithmetic is worse but
because the quantity does not exist in its vocabulary.

% =====================================================================================
\section{The Bayesian model}
\label{sec:its:bayes}

The Bayesian arm fits the \emph{pre-period only} and extrapolates, which is the more literal
translation of the ITS idea than \cref{eq:its:seg} is. With $x_t = (1,\ t,\ \sin(2\pi t/7),\
\cos(2\pi t/7))^{\top}$ the design row, the model as implemented by CausalPy
\citep{causalpy} on PyMC \citep{salvatier2016} is
%
\begin{align}
  Y_t \mid \beta, \sigma &\;\sim\; \mathcal N\big(x_t^{\top}\beta,\ \sigma^{2}\big),
      \qquad t < t^{*}, \label{eq:its:lik}\\
  \beta_j &\;\sim\; \mathcal N\big(0,\ \nbTenPriorSd^{2}\big), \qquad j = 0,\dots,3,
      \label{eq:its:prior_b}\\
  \sigma &\;\sim\; \text{HalfNormal}(1). \label{eq:its:prior_s}
\end{align}
%
One sentence of justification is owed whenever defaults are used. On a conversion scale where the
level is $\approx 0.10$ and day-to-day noise is $\approx 0.006$, a prior standard deviation of
\nbTenPriorSd{} on every coefficient and a HalfNormal$(1)$ on $\sigma$ are \emph{extremely}
diffuse --- effectively flat --- so the \nbTenNPre{} pre-period days dominate and the posterior is
likelihood-driven. As throughout this book, the prior is doing very little work; the likelihood is
doing all of it. (On a differently scaled outcome --- raw daily revenue in euros --- the same
defaults would distort or slow the fit, and one would rescale the outcome or tighten the priors.
A default is only defensible once it has been checked against the units.)

The counterfactual for a post-redesign day is the posterior predictive of \cref{eq:its:lik}
evaluated at $t \ge t^{*}$, and the effect is the posterior of the pointwise gap:
%
\begin{equation}
  \hat\Delta_t \;=\; Y_t - \widehat{Y_t(0)},
  \qquad \widehat{Y_t(0)} \;\sim\; p\big(\tilde Y_t \mid \mathbf y_{<t^{*}}\big),
  \qquad t \ge t^{*} .
  \label{eq:its:gap}
\end{equation}
%
Each draw of $(\beta, \sigma)$ yields a whole counterfactual path, hence a whole path of gaps,
hence posterior \emph{distributions} for both $\bar\Delta$ and $C_T$ rather than single numbers.
The band widens with the horizon --- visibly, in \cref{fig:its:its} --- which is the model saying
the one honest thing it can say about extrapolation: a counterfactual grows less certain the
further it is projected. It is worth being precise about what that widening does and does not
represent. It is uncertainty about the \emph{fitted line}, propagated forward. It is not
uncertainty about whether the line was the right one to project.

Inference is by the No-U-Turn sampler \citep{hoffman2014}: \nbTenNDraws{} post-warmup draws, with
$\hat R = \nbTenRhat$, a minimum bulk effective sample size of \nbTenEss{} \citep{vehtari2021},
and \nbTenDivergences{} divergent transitions. The chains converged; nothing below is a sampler
artefact.

The posterior mean lift is \pp{\nbTenBayesAvg}, with a 90\,\% credible interval of
$[\pp{\nbTenBayesLo},\ \pp{\nbTenBayesHi}]$, against the planted \pp{+\nbTenTrueLift}
(\cref{fig:its:its}). The cumulative gap over the window is $C_T = \nbTenCumBayes$ pp-days, which
\cref{sec:its:euro} converts to euros.

\input{build/floats/nb10_its}

% =====================================================================================
\section{Diagnostics and validation}
\label{sec:its:valid}

Five checks, and the most important thing about them is stated at the end of the section rather
than the beginning: \emph{every one of them interrogates the pre-period model}, and the
identifying assumption lives in the post-period.

\subsection*{Placebo in time}

\Cref{as:its:antic} is tested by moving the intervention backwards. The whole method is refit with
a fake redesign on day \nbTenFakeDay, inside the true pre-period, where by construction nothing
happened. It must return approximately zero, and does: \pp{\nbTenPlacebo}. A sweep over further
fake dates (\nbTenPlaceboDates) returns at most \pp{\nbTenPlaceboMax} in absolute value
(\cref{fig:its:diag}, right). The method does not manufacture lift where there is none.

\subsection*{Residual autocorrelation, and what an alarm looks like}

\Cref{as:its:form} and \cref{as:its:dep} are both tested by the autocorrelation of the pre-period
fit's residuals. With the weekly Fourier pair in the formula, \nbTenAcfOut{} of ten lags fall
outside the $\pm 2/\sqrt{n}$ band of $\pm\nbTenAcfBand$; the lag-1 bar is \nbTenAcfOne{} and the
lag-7 bar \nbTenAcfSeven. The residuals are effectively white, so the intervals are not understated.

A green diagnostic teaches nothing about what an alarm looks like, so the instrument is fired on
purpose. Refit the \emph{same data} with the seasonal pair deleted --- \code{conversion \textasciitilde{}
1 + t} --- and the lag-7 residual bar jumps to \nbTenAcfBadSeven, several times the band
(\cref{fig:its:diag}, left). That is the textbook signature of unmodelled weekly seasonality, and
it is the reason the ACF panel is not decoration. The damage is instructive: the unmodelled sine
lands in the residual, so $\hat\sigma$ inflates from \pp{\nbTenResidSd} to \pp{\nbTenResidSdBad}
($\nbTenResidSdRatio\times$) and the credible interval widens with it, from \pp{\nbTenOkWidth} to
\pp{\nbTenBadWidth} ($\nbTenBadWidthRatio\times$), while the point estimate barely moves
(\pp{\nbTenBadAvg}) --- only because the weekly cycle happens to nearly average out over a
post-window of \nbTenNPost{} days, which is a whole number of weeks. With a ragged window, or with
seasonality larger relative to the effect, the point estimate moves too. When the lag-7 bar jumps
the band, fix the formula before trusting either the number or the band.

\input{build/floats/nb10_diag}

\subsection*{Earning the extrapolation: a backtest}

The whole method is a forecast, so the forecast is rehearsed where it can still be graded. Hold out
the last \nbTenHold{} pre-period days, refit on days $0$ to $\nbTenRedesignDay - \nbTenHold - 1$
only, and extrapolate across the holdout exactly as the real fit extrapolates past $t^{*}$. The
90\,\% posterior-predictive band covers \nbTenBtCovered{} of \nbTenHold{} held-out days
(\pc{\nbTenBtCovPct}, against a nominal 90) with a mean forecast error of \pp{\nbTenBtErr}
(\cref{fig:its:backtest}, left). The level--trend--weekly shape extrapolates \nbTenHold{} days
ahead without drifting.

This is the most direct check ITS possesses --- and it is exactly a \emph{rehearsal}. The day-$t^{*}$
counterfactual is the same forecast pushed \nbTenNPost{} days instead of \nbTenHold, into a window
where no grading is possible, and across a date on which something changed by hypothesis. A
successful backtest earns the extrapolation. It does not earn \cref{as:its:shock}.

\subsection*{How much history is enough}

Refitting on the last $L$ pre-period days only shows what the design lives or dies on
(\cref{fig:its:backtest}, right). Shortening the history does not \emph{bias} the estimate here ---
the point stays near the truth at every $L$ --- it \emph{blinds} it: the credible interval at
$L = \nbTenLenShort$ days is \pp{\nbTenLenShortWidth} wide against \pp{\nbTenLenLongWidth} at
$L = \nbTenLenLong$, a factor of $\nbTenLenWidthRatio\times$. With roughly four weekly cycles, the
trend and the step blur into each other and neither is identified. On real data a short window can
bias outright as well, if it happens to sit on a local anomaly --- \cref{as:its:date}'s trap. A
long, boring pre-period is a feature, not a cost.

\input{build/floats/nb10_backtest}

\subsection*{Recovery across fresh series}

One sample can land lucky. Across \nbTenNSeed{} fresh series from \cref{eq:its:dgp}, the estimator
recovers a mean of \pp{\nbTenSeedMean} against the planted \pp{+\nbTenTrueLift} --- a bias of
\pp{\nbTenSeedBias}, with a seed-to-seed standard deviation of \pp{\nbTenSeedSd} --- and the
90\,\% credible interval covers the truth on \nbTenSeedCov{} of \nbTenNSeed{} series
(\pc{\nbTenSeedCovPct}). The estimator is centred, and its interval is consistent with being
calibrated. (With this many trials the tally is binomial noise around the nominal count, not a
precise coverage estimate; it should be read as ``consistent with calibrated'', which is all a
handful of refits can establish.)

\subsection*{The scope of every check above}

Placebo, ACF, backtest, history sweep, multi-seed recovery: all five are green, and all five
interrogate \textbf{the pre-period model}. Not one of them looks at the post-period, because there
is nothing in the post-period to look at --- the counterfactual is missing, which is why the method
exists. \Cref{sec:its:wrong} exploits precisely that.

% =====================================================================================
\section{Point estimate versus posterior}
\label{sec:its:compare}

\input{build/tables/nb10_compare}

\subsection*{On the number, they agree --- and on the width, neither wins}

\Cref{tab:its:compare} puts both arms on the estimand of \cref{eq:its:estimand}. The point
estimates differ by \pp{\nbTenEstDiff} --- \pc{\nbTenEstDiffPct} of the effect --- and both
intervals cover the planted \pp{+\nbTenTrueLift}. This is not a coincidence and not a
disappointment: \textbf{the causal work was done by the identification, not by the machinery.}
Once \cref{eq:its:ident} is granted, a Bayesian regression with diffuse priors and \nbTenNPre{}
days of data is doing arithmetic that five OLS regressors had already done. The name for this is
the Bernstein--von Mises phenomenon \citep{gelman2013}: with a well-specified likelihood and enough
data, the posterior concentrates where the maximum-likelihood estimate does.

The agreement is the more impressive for the two arms not being the same fit. The Bayesian ITS fits
the pre-period \emph{only} and extrapolates; the segmented regression fits the \emph{whole} series
and lets $\delta$ and $\gamma$ absorb the intervention --- so it borrows the post-period days to
pin down the trend and the weekly cycle, buying efficiency at the price of letting post-period
structure shape the counterfactual. Estimating the extra slope parameter $\gamma$ spends that
efficiency again. The two half-widths come out at \pp{\nbTenClHalf} (classical) and
\pp{\nbTenBayesHalf} (Bayesian), a ratio of $\nbTenHalfRatio\times$, which is far inside noise.
\textbf{Neither arm is reliably the sharper}, and this chapter declines to call a winner it cannot
demonstrate.

\subsection*{The mirage: the cumulative-impact interval was not the payoff}

The obvious pitch for Bayes here is the \emph{derived} quantity. \Cref{sec:its:euro} needs the
realized 60-day euro figure with an interval attached, and the posterior hands it over: push every
draw of $C_T$ through the euro conversion and read off the quantiles. It looks like exactly the
sort of thing a posterior is for --- and it is a mirage.

Look at \cref{eq:its:cum}. $C_T = \nbTenNPost \times \bar\Delta$ is a \textbf{linear} function of
the very coefficient the classical arm already estimated, and \cref{sec:its:euro} then multiplies
by a \emph{constant} (\nbTenVisitors{} visitors $\times$ \eur{\nbTenConvValue} per conversion).
Scale the classical HAC interval by those constants and it reproduces the posterior almost exactly
(\cref{tab:its:derived}): \eur{\nbTenEurCls} classically against \eur{\nbTenEurBayes} Bayesianly
--- a gap of \eur{\nbTenEurGap} on a figure of half a million --- with interval widths in a ratio
of $\nbTenEurWidthRatio\times$. No delta method. No bootstrap. No gymnastics.

\input{build/tables/nb10_derived}

Where propagating uncertainty through a \emph{linear} map is all that is asked, the classical arm
propagates it \emph{by multiplication}, and Bayes buys convenience, not capability. The general
statement is worth extracting, because it is the discipline that keeps this book's comparisons
honest: a posterior's ability to push draws through an arbitrary function is a real advantage
exactly when the function is \emph{non}-linear, or when several uncertain quantities are combined.
It is worth nothing at all when the function is a multiplication by a known constant. Before
claiming a derived quantity as a Bayesian payoff, check which of the two it is.

\subsection*{What the posterior did buy}

One thing, and it is the thing the business asked for. The last column of \cref{tab:its:compare}
is $\Prob(\bar\Delta > c)$, and the classical entry reads \textbf{not defined}. That is the literal
truth, not a rhetorical concession: a confidence interval is a statement about a procedure and a
$p$-value ranks data against a null; neither attaches a probability to a hypothesis about a fixed
parameter. There is no amount of effort that converts one into the other.

And \cref{sec:its:euro}'s decision is nothing \emph{but} such a probability. It converts the
\eur{\nbTenCost} build cost into a break-even lift and asks $\Prob(\text{annual value} >
\text{cost})$ once per persistence scenario --- a whole distribution travelling through a
non-linear decision rule, four times, emerging as four probabilities that flip a verdict between
two rows. No confidence interval can be tortured into producing that table. \emph{That} is the
non-linear map the previous subsection said to look for, and it is where the posterior finally
earns its place.

Bayes did not rescue the estimate here; the estimate was never in danger. It bought the right to
make a probability statement about the effect.

It is worth being explicit about what did \emph{not} go wrong, because a neighbouring chapter of
this book turns on exactly that failure. In the synthetic-control chapter the Bayesian interval
under-covered badly, and the diagnosis was a \emph{misspecified likelihood} --- an iid error term
laid over a random walk --- for which \citet{mueller2013} supplies the principled repair: rescale
the posterior by a sandwich covariance and its frequentist coverage returns. None of that applies
here. The residuals of \cref{eq:its:lik} really are close to white (\cref{sec:its:valid}), the
posterior's scale is right, and its coverage across fresh series is \pc{\nbTenSeedCovPct}. The
Bayesian arm of this chapter is not overconfident about \emph{sampling noise}. It is
overconfident about nothing at all --- and it still gets \cref{sec:its:wrong} completely wrong,
which is the point.

\subsection*{And what nothing bought}

The ledger, stated in both directions. Bayes did \textbf{not} buy a better point estimate (they
agree to \pp{\nbTenEstDiff}). It did \textbf{not} buy a tighter interval (the half-widths agree to
$\nbTenHalfRatio\times$, with neither arm reliably sharper). It did \textbf{not} buy the
realized-euro interval (a linear rescaling of the classical one). And, as the next section shows,
it did \textbf{not} buy identification --- which is the only one of the four that would have
mattered.

% =====================================================================================
\section{The decision, in euros}
\label{sec:its:euro}

The site takes \nbTenVisitors{} visitors a day, a conversion is worth \eur{\nbTenConvValue}, and
the redesign's amortized build-and-maintenance cost is \eur{\nbTenCost} a year. Two kinds of euro
follow from those numbers and they must not be mixed.

\subsection*{Realized: a fact about the window that was observed}

The cumulative posterior $C_T = \nbTenCumBayes$ pp-days converts directly into
\nbTenExtraConv{} extra conversions over the \nbTenNPost{} observed days, or
\eur{\nbTenEurBayes} of value, 90\,\% interval $[\eur{\nbTenEurBayesLo},\ \eur{\nbTenEurBayesHi}]$
--- \pc{\nbTenEurRealizedPctCost} of the \eur{\nbTenCost} build cost, \emph{already banked},
with no persistence assumption invoked. (And, per \cref{sec:its:compare}, computable to within
\eur{\nbTenEurGap} from the classical arm alone.)

\subsection*{Projected: a bet on a parameter the data cannot identify}

A redesign lift is rarely permanent --- \textbf{novelty fades}. Projecting a 60-day lift flat
across a year contradicts this chapter's own message about extrapolation, so the projection is made
against an explicit decay. If the effect decays with a novelty half-life $h$, the year is worth
%
\begin{equation}
  \text{effective days} \;=\; \sum_{t=0}^{364} 0.5^{\,t/h},
  \qquad
  \text{annual value} \;=\; \bar\Delta \cdot \text{\nbTenVisitors} \cdot \text{\eur{\nbTenConvValue}}
                            \cdot \text{effective days},
  \label{eq:its:decay}
\end{equation}
%
and the decision rule --- the convention used throughout this book --- is applied to the posterior:
%
\begin{equation}
  \Prob(\text{value} > \text{cost})\ \begin{cases}
    > 0.9 & \Rightarrow\ \textsc{keep},\\
    < 0.5 & \Rightarrow\ \textsc{roll back / renew},\\
    \text{otherwise} & \Rightarrow\ \textsc{renew and re-test}.
  \end{cases}
  \label{eq:its:rule}
\end{equation}

The break-even lift is \pp{\nbTenBeLift} --- the lift at which \nbTenVisitors{} visitors a day, at
\eur{\nbTenConvValue} a conversion, over \nbTenHorizon{} days, exactly cover the \eur{\nbTenCost}
build --- and the measured lift is $\nbTenLiftOverBe\times$ it. \Cref{tab:its:persistence} sweeps
the half-life.

\input{build/tables/nb10_persistence}
\input{build/floats/nb10_euro}

The verdict is \textbf{\nbTenPersVerdictnodecay} under every scenario except one, and the exception
is the instructive row: under a one-month novelty burn-out the lift is gone before it has paid, and
$\Prob(\text{value} > \text{cost}) = \nbTenPersPonem$ --- \nbTenPersVerdictonem. Notice what has
happened to the decision. It no longer hinges on \emph{whether} the redesign worked
($\Prob(\bar\Delta > 0) = \nbTenPPos$) but on \emph{how long the lift persists} --- a parameter a
60-day posterior structurally cannot identify, because decay reveals itself only over a longer
horizon. The \textsc{keep} call survives down to a half-life of about \nbTenHalflifeNinety{} days
and dies below \nbTenHalflifeFifty{} (\cref{fig:its:euro}, right).

The honest deliverable is therefore not an annualised number. It is the persistence table plus a
\textbf{monitoring plan}: re-estimate at day 120 and day 180 and see which row of
\cref{tab:its:persistence} the business is living in.

\subsection*{Headroom: what to do with an untestable assumption}

\Cref{as:its:shock} cannot be tested. It can, however, be \emph{priced}, and pricing it is the
honest replacement for pretending otherwise. Ask: how large would a coincident shock have to be
before the \textsc{keep} call flips? Even at the conservative \pp{+\nbTenLiftLoNinety} lower bound
of the 90\,\% credible interval, a concurrent shock would have to account for more than
\pp{\nbTenHeadroomPp} of the measured lift --- the distance to the \pp{\nbTenBeLift} break-even ---
before the no-decay annual value falls to cost. That is wide headroom, and it is a sentence a CMO
can act on, in a way that ``we assume no concurrent shock'' is not.

The memo writes itself. The redesign works, and has already banked \pc{\nbTenEurRealizedPctCost} of
its cost in realized conversion value. The open question is how long the lift lasts, not whether it
exists. And the estimate is only as good as the launch calendar --- which is the subject of the
next section.

% =====================================================================================
\section{What can go wrong}
\label{sec:its:wrong}

\subsection*{A concurrent shock --- and the identical failure of both arms}

Every diagnostic in \cref{sec:its:valid} came back green. Every one of them tested the pre-period
model. \Cref{as:its:shock} lives in the post-period, so the only honest way to teach it is to break
it on purpose.

Take the same series, the same \pp{+\nbTenTrueLift} redesign, and add an \emph{unrelated}
\pp{+\nbTenShock} shock arriving on the redesign date --- a competitor's outage, a press mention, a
coincident price cut. The redesign is still worth \pp{+\nbTenTrueLift}. What do the estimators
report?

\input{build/tables/nb10_shocktab}
\input{build/floats/nb10_shockfig}

Both of them report the sum, \pp{\nbTenShockSum}, and both report it \emph{confidently}. The
Bayesian ITS returns \pp{\nbTenShockBayes} with a 90\,\% credible interval of
$[\pp{\nbTenShockBayesLo},\ \pp{\nbTenShockBayesHi}]$, which excludes the true
\pp{+\nbTenTrueLift}. The classical segmented regression, refit on the very same series with its
Newey--West sandwich intact, returns \pp{\nbTenShockCls} with a 90\,\% confidence interval of
$[\pp{\nbTenShockClsLo},\ \pp{\nbTenShockClsHi}]$, which excludes the true \pp{+\nbTenTrueLift}
just as tightly. \textbf{The same error, of the same size, with the same confidence.}

And nothing flinched. The sampler converged ($\hat R = \nbTenShockRhat$). The pre-period fit is as
good as it ever was. The placebo at day \nbTenFakeDay{} \emph{still passes}, returning
\pp{\nbTenShockPlacebo}, because the pre-period is untouched by a post-period shock --- the
falsification test is not fooled so much as looking elsewhere.

That is the precise scope of each tool, and it deserves to be stated as a definition rather than a
caveat:

\begin{itemize}[leftmargin=*]
  \item the \textbf{placebo} validates the \emph{pre-period model} --- ``does the method invent
        lift where none exists?'' --- and says nothing about post-period attribution;
  \item the \textbf{credible interval} and the \textbf{confidence interval} both quantify
        \emph{sampling noise} --- ``how precisely do we know the gap between what happened and what
        the model projected?'' --- and neither quantifies \emph{identification error} --- ``is that
        gap all redesign?''
\end{itemize}

\begin{remark}[The moral, stated precisely]
The posterior is not more honest about identification failure than the confidence interval is. It
is \emph{exactly} as confident, and \emph{exactly} as wrong. The credible interval and the
confidence interval both quantify sampling noise, and \textbf{the ITS problem is not sampling
noise} --- it is that the post-period counterfactual is unobservable and the assumption that
reconstructs it is untestable. No sandwich estimator defends \cref{as:its:shock}; no prior defends
it; no sampler defends it. \textbf{The fix is never a better estimator. It is a better design.}
\end{remark}

The design fixes, in order of preference. \emph{Stagger the rollout} across stores, cohorts or
regions, so that at any moment some units are untreated and difference-in-differences applies
(the staggered-rollout chapter) --- this is nearly always available for a website, and refusing
it is a choice.
\emph{Hold out a geography} and reconstruct the counterfactual from live untreated units rather
than from history, as the synthetic-control chapter does. Failing both, \emph{audit the launch
calendar by hand} --- pricing, campaigns, competitor moves, PR, tracking changes, in a window
around the date --- and then \emph{quantify the headroom}, as \cref{sec:its:euro} does. What is not
on the list is a better estimator. \Cref{tab:its:shock} is the proof that it would not have helped.

There is one thing a Bayesian treatment can genuinely add here, and it is worth naming so that the
lesson is not over-read into ``Bayes is useless for ITS''. \Citet{brodersen2015}'s Bayesian
structural time-series approach --- \code{CausalImpact} --- models the counterfactual with a local
level, a local trend and, crucially, \emph{contemporaneous covariates}: other series that respond
to the same shocks but not to the treatment. A covariate that moves with the competitor's outage
but not with the redesign would absorb the shock and rescue the estimate. That is a genuine repair
--- but note what it is. It is not a better \emph{inference}; it is the smuggling of a control
group back into a problem that did not have one. The fix is still design. It has simply been
disguised as a model.

\subsection*{Regression to the mean --- the subtle version of the same failure}

\Cref{as:its:date} fails quietly and often. Redesigns ship because conversion \emph{looked bad};
if a transient dip triggered the launch, the pre-period fit is dragged down by the dip and the
natural rebound is booked as redesign lift. This is \cref{as:its:shock} in another costume --- the
``other term'' at $t^{*}$ is the mean-reversion of the metric itself --- and it is more insidious
because it requires no external event at all. Always ask \emph{why this date}. If the metric
triggered the launch, exclude or explicitly model the dip window, and expect the placebo-date sweep
of \cref{fig:its:diag} to light up near it.

\subsection*{A misspecified pre-period model}

\Cref{as:its:form} fails loudly, which makes it the least dangerous of the three. Deleting the
weekly cycle from the formula sends the lag-7 residual ACF to \nbTenAcfBadSeven{} and inflates the
credible interval by $\nbTenBadWidthRatio\times$ (\cref{sec:its:valid}). The instrument works. The
point estimate survived here only because the post-window happened to be a whole number of weeks;
do not count on that twice.

\subsection*{Persistent residuals}

\Cref{as:its:dep} fails silently under an iid standard error, whose coverage falls to
\pc{\nbTenArCovIidEight} at $\rho = 0.8$ (\cref{tab:its:hac}) while its reported precision does not
move at all. HAC repairs most of it, to \pc{\nbTenArCovHacEight}, and structure --- an explicit AR
term, or a state-space model --- repairs more. The relevant discipline is simply that the
covariance assumption must be \emph{made}, out loud, and never accepted by default.

% =====================================================================================
\section{Takeaways}
\label{sec:its:take}

\begin{enumerate}[leftmargin=*, itemsep=.45em]
  \item \textbf{The past is the control group, and it is the weakest control group there is.} When
        everyone is treated at once, the counterfactual can only be extrapolated from history. ITS
        does that honestly --- it strips \pp{\nbTenClDriftRemoved} of pre-existing drift out of the
        dashboard's \pp{\nbTenNaive} and lands on \pp{\nbTenClAvg} against a planted
        \pp{+\nbTenTrueLift} --- but the counterfactual is a forecast into a window where nothing
        can check it.

  \item \textbf{The covariance choice is an argument, not a default.} On this series Newey--West
        and iid agree to $\nbTenClSeRatio\times$, because the noise is clean and the weekly cycle is
        in the formula. Rebuild the same series with persistent errors and the iid interval's
        coverage collapses from \pc{\nbTenArCovIidZero} to \pc{\nbTenArCovIidEight} while its
        reported standard error does not move --- it is not underestimating the uncertainty, it is
        structurally unable to perceive it. HAC holds \pc{\nbTenArCovHacEight}. It is insurance that
        costs nothing when the residuals are clean, and one cannot know which world one is in
        without computing it.

  \item \textbf{Classical and Bayesian agree, on the number and on the width.} They differ by
        \pp{\nbTenEstDiff} on the estimate and by $\nbTenHalfRatio\times$ on the half-width. Neither
        arm is reliably the sharper, and the chapter refuses to call a winner it cannot demonstrate.

  \item \textbf{The cumulative-euro interval was not a Bayesian payoff --- it was a linear
        rescaling.} $C_T = \nbTenNPost\,\bar\Delta$ times a constant is a linear map of a
        coefficient the classical arm already had, and scaling the HAC interval reproduces the
        posterior to \eur{\nbTenEurGap}. Bayes buys convenience there, not capability. Reserve the
        claim of a Bayesian payoff for genuinely \emph{non}-linear derived quantities --- of which
        \cref{tab:its:persistence} is one.

  \item \textbf{What the posterior did buy is $\Prob(\bar\Delta > c)$.} The classical column reads
        \emph{not defined}, literally. The euro rule of \cref{eq:its:rule} consumes nothing else,
        and the persistence table is a whole distribution pushed through a non-linear decision rule
        four times. Use the posterior for the decision, and design-based reasoning for your
        confidence in it.

  \item \textbf{Nothing bought protection from identification failure --- and this is the chapter.}
        An unrelated \pp{+\nbTenShock} shock on the redesign date makes both arms report
        \pp{\nbTenShockSum}, each inside a tight interval that excludes the truth, while every
        in-sample diagnostic stays green and the placebo still passes. The credible interval and the
        confidence interval both quantify \emph{sampling noise}, and the ITS problem is not sampling
        noise: the post-period counterfactual is unobservable and the assumption that reconstructs
        it is untestable. \textbf{The fix is never a better estimator; it is a better design} ---
        stagger the rollout so difference-in-differences applies, keep a holdout
        geography for synthetic control, or, failing both, audit the launch calendar by hand and
        quantify
        the headroom.

  \item \textbf{The decision stopped being about the effect and started being about its
        half-life.} At $\Prob(\bar\Delta > 0) = \nbTenPPos$ and a lift
        $\nbTenLiftOverBe\times$ break-even, \textsc{keep} survives every persistence scenario but a
        one-month burn-out. The deliverable is therefore \cref{tab:its:persistence} plus a
        monitoring plan --- not an annualised number that quietly assumes the answer.
\end{enumerate}

\notebooknote{%
  The analysis of \crefrange{sec:its:dgm}{sec:its:wrong} is
  \code{notebooks/10\_redesign\_its.ipynb}, which runs in the legacy environment
  (\code{pymc<6}, kernel \code{cmp-legacy}) because it uses CausalPy's
  \code{InterruptedTimeSeries}. It has a fast teaching mode (\code{CMP\_FAST=1}, short chains) and
  a full mode (\code{CMP\_FAST=0}); every number in this chapter comes from the full run, emitted by
  \code{cmp.report} and injected here as a macro. The simulator is \code{cmp.dgp.its\_redesign}, the
  segmented regression is \code{cmp.classical.segmented\_its}, and the seeds are fixed in the
  notebook. The canonical public ITS exercise remains the effect of anti-smoking legislation on
  cigarette sales; \citet{bernal2017} is the tutorial to hand a new analyst, and
  \citet{brodersen2015} the model to reach for when contemporaneous control series exist.}
